\chapter{From Cross-Entropy to Mean-Neutral Attention: A First-Principles Derivation}
\label{chap:laplacian_softmax}

\section{Introduction}

The Transformer architecture is built upon a simple but powerful principle:
the meaning of a token is determined by its context.
This principle is implemented through the self-attention mechanism,
which aggregates information across tokens via a softmax-weighted average.

Empirically, however, deep Transformer networks exhibit strong anisotropy.
Hidden representations accumulate a coherent mean bias and develop
dominant spectral directions.
This phenomenon persists even when the loss function itself
does not explicitly encourage such bias.

In this chapter we ask a precise mathematical question:

\emph{Does the standard softmax attention operator inherently generate mean drift?}

We begin from the training objective itself, examine the structure of
cross-entropy gradients, and derive the implications for contextual operators.
We then show that softmax attention behaves as a low-pass averaging operator,
which naturally produces coherent drift.
Finally, we derive an alternative operator---the Laplacian softmax---
from first principles of contrastive learning,
and show that it is intrinsically mean-neutral.

This derivation does not rely on empirical observation.
It follows directly from the mathematics of the objective.

\section{Cross-Entropy Training and Translation Invariance}

Large language models are trained with next-token prediction.
Let $h_t \in \mathbb{R}^H$ denote the hidden representation at position $t$,
and let $W_U$ be the output projection.
The logits are

\begin{equation}
z_t = h_t W_U.
\end{equation}

The probability of the next token is

\begin{equation}
p_t = \mathrm{softmax}(z_t).
\end{equation}

The loss is cross-entropy:

\begin{equation}
\mathcal{L}
=
-\sum_{t} \log p_t(y_t).
\end{equation}

Consider shifting all hidden states by a constant vector $c$:

\begin{equation}
h_t' = h_t + c.
\end{equation}

The logits become

\begin{equation}
z_t' = z_t + c W_U.
\end{equation}

If $cW_U$ is constant across vocabulary indices,
softmax remains unchanged:

\begin{equation}
\mathrm{softmax}(z_t + \alpha \mathbf{1}) = \mathrm{softmax}(z_t).
\end{equation}

Therefore, the loss is invariant to global hidden-state shifts.

\paragraph{Implication.}
Cross-entropy does not penalize coherent mean drift.
Nothing in the objective enforces $\mu(h_t) = 0$.

\section{Gradient Structure and Contrastive Learning}

The gradient of the loss with respect to $h_t$ is

\begin{equation}
\nabla_{h_t} \mathcal{L}
=
W_U^\top (p_t - y_t).
\end{equation}

Observe that this gradient depends only on the difference
between predicted distribution $p_t$ and target $y_t$.
It does not depend on absolute magnitude of $h_t$.

Thus training is inherently \emph{contrastive}:
we push correct logits up and incorrect logits down.
Learning is driven by relative differences,
not absolute scale.

This is the first principle we must respect:
semantic learning is fundamentally contrastive.

\section{Softmax Attention as Convex Averaging}

Standard self-attention is defined as

\begin{equation}
Q = XW_Q, \quad
K = XW_K, \quad
V = XW_V,
\end{equation}

\begin{equation}
A = \mathrm{softmax}\!\left(\frac{QK^\top}{\sqrt{d}}\right),
\end{equation}

\begin{equation}
Y = AV.
\end{equation}

The softmax is applied row-wise, so

\begin{equation}
A_{ij} \ge 0,
\quad
\sum_j A_{ij} = 1.
\end{equation}

Each output token is therefore a convex combination:

\begin{equation}
y_i = \sum_j A_{ij} v_j.
\end{equation}

This has a crucial geometric consequence.

\subsection{Mean Preservation Property}

Let

\begin{equation}
V = \mathbf{1}\mu^\top + V_c,
\end{equation}

where $\mu$ is the shared mean component
and $V_c$ is zero-mean across tokens.

Then

\begin{equation}
Y = A(\mathbf{1}\mu^\top + V_c)
=
\mathbf{1}\mu^\top + AV_c.
\end{equation}

Because $A\mathbf{1} = \mathbf{1}$,
the shared component propagates unchanged.

Thus softmax attention preserves and spreads coherent mean.

\section{Why Softmax Generates Mean Drift}

Softmax weights are strictly positive.
They enforce cooperative aggregation.
This is a low-pass filter:
shared components survive,
differences are averaged.

When combined with residual connections,

\begin{equation}
X_{\ell+1} = X_\ell + AV_\ell,
\end{equation}

any small coherent component
accumulates layer by layer.

Thus even if gradients are contrastive,
the forward operator is aggregative.
This mismatch explains mean drift.

\section{Contrastive Principle for Contextual Operators}

The semantic assumption underlying attention is:

\begin{quote}
A token should move toward relevant context and away from irrelevant context.
\end{quote}

Standard softmax only implements attraction.
It does not implement repulsion.
Every token contributes positively.

However, semantic structure is bipolar:
tokens are correlated in some subspaces
and anticorrelated in others.

Therefore a more faithful operator should implement
difference-based updates.

\section{Deriving Laplacian Softmax from First Principles}

We seek an operator that:

\begin{enumerate}
\item Preserves contrastive learning logic.
\item Removes shared offsets.
\item Maintains contextual selectivity.
\end{enumerate}

Define the difference attention update:

\begin{equation}
\Delta y_i
=
\sum_j A_{ij}(v_j - v_i).
\end{equation}

Expanding:

\begin{equation}
\Delta y_i
=
\left(\sum_j A_{ij} v_j\right) - v_i.
\end{equation}

Summing across tokens:

\begin{equation}
\sum_i \Delta y_i
=
\sum_i \sum_j A_{ij} v_j
-
\sum_i v_i.
\end{equation}

Since rows of $A$ sum to one,

\begin{equation}
\sum_i \sum_j A_{ij} v_j
=
\sum_j v_j.
\end{equation}

Therefore,

\begin{equation}
\sum_i \Delta y_i = 0.
\end{equation}

Thus

\begin{equation}
\mu(\Delta Y) = 0.
\end{equation}

This operator is mean-neutral.

\subsection{Laplacian Interpretation}

Define graph Laplacian:

\begin{equation}
L = I - A.
\end{equation}

Then

\begin{equation}
\Delta Y = -L V.
\end{equation}

This is a diffusion operator.
It removes constant modes.
Constant vectors lie in the nullspace of $L$.

Thus Laplacian attention explicitly removes shared components.

\section{Geometric Consequences}

Softmax attention:
\begin{equation}
Y = AV
\end{equation}
is a convex smoothing operator.

Laplacian attention:
\begin{equation}
Y = V - LV
\end{equation}
is contrastive and mean-neutral.

Softmax:
Low-pass filtering.

Laplacian:
High-pass filtering.

In high dimension, even small coherent bias scales as
\begin{equation}
\|\mu\| \sim \sqrt{H}\epsilon.
\end{equation}

Thus preventing mean regeneration is essential.

\section{Conclusion}

From the cross-entropy objective,
we derive that learning is inherently contrastive.
However, standard softmax attention performs only
positive aggregation.
This mismatch creates coherent drift.

By re-deriving attention from a contrastive principle,
we obtain Laplacian softmax,
which is provably mean-neutral
and suppresses shared bias modes.

This operator aligns the geometry of contextual learning
with the mathematics of the objective,
and provides a principled foundation
for mean-stable Transformer architectures.


\chapter{Replacing SwiGLU: Mean-Neutral Nonlinearities for Transformer Architectures}

\section{Introduction}

Modern Transformer architectures rely heavily on gated feed-forward
nonlinearities, most prominently SwiGLU.
Empirically, SwiGLU improves training stability and performance compared to
ReLU or GELU in large language models.
However, from a geometric perspective, SwiGLU introduces a systematic mean
shift in activations, even when the input is zero-mean.
In deep residual architectures, this mean shift accumulates layer by layer,
producing coherent bias growth and contributing to anisotropy.

This chapter develops a first-principles analysis of why SwiGLU regenerates
mean bias, how this affects representation geometry, and how one can design
mean-neutral alternatives without sacrificing expressive power.
Our goal is not merely to eliminate bias, but to design nonlinear operators
whose geometry is aligned with the high-dimensional statistical structure of
language.

\section{Mean Regeneration in Two-Layer Feed-Forward Networks}

Consider a standard two-layer Transformer MLP block without residuals:

\begin{equation}
Z = X W_1,
\qquad
A = \phi(Z),
\qquad
Y = A W_2,
\end{equation}

where $X \in \mathbb{R}^{N \times H}$ contains token representations,
$W_1 \in \mathbb{R}^{H \times M}$,
$W_2 \in \mathbb{R}^{M \times H}$,
and $\phi$ is a nonlinear activation.

Assume that $X$ is zero-mean across tokens:

\[
\mu(X) = \frac{1}{N} \mathbf{1}^\top X = 0.
\]

Even under this assumption, the nonlinearity $\phi$ may generate
a nonzero mean:

\[
\mu(A) = \frac{1}{N} \mathbf{1}^\top \phi(XW_1).
\]

If $\phi$ is not an odd function,
then even when $Z$ is symmetrically distributed,
$\mathbb{E}[\phi(Z)] \neq 0$.
This creates a coherent drift vector.

For ReLU:
\[
\phi(z) = \max(0,z),
\]
we have
\[
\mathbb{E}[\phi(z)] > 0
\quad \text{whenever } z \sim \mathcal{N}(0,\sigma^2).
\]

For GELU or SiLU (used inside SwiGLU), the same holds:
\[
\mathbb{E}[\phi(z)] > 0.
\]

Thus the nonlinearity alone introduces mean shift.

After multiplication by $W_2$, this mean becomes dense in output space:

\[
\mu(Y) = \mu(A) W_2.
\]

Hence even if $X$ is centered,
$Y$ is not.

In residual networks:

\[
X_{\ell+1} = X_\ell + Y_\ell,
\]

the mean accumulates additively:

\[
\mu(X_{\ell+1}) = \mu(X_\ell) + \mu(Y_\ell).
\]

Therefore, non-odd activations cause systematic mean growth.

\section{SwiGLU and Mean Amplification}

SwiGLU is defined as

\begin{equation}
\text{SwiGLU}(X)
=
(X W_1) \odot \text{SiLU}(X W_2),
\end{equation}

where $\odot$ denotes elementwise multiplication.

Even if $X$ is zero-mean,
both $XW_1$ and $XW_2$ are approximately zero-mean (by linearity).
However, SiLU is non-odd, so

\[
\mu(\text{SiLU}(XW_2)) > 0.
\]

The product of a symmetric random variable and a positively skewed one
is also skewed.
Thus SwiGLU compounds mean generation through multiplicative interaction.

Empirically, repeated SwiGLU layers increase
the mean bias ratio

\[
R = \frac{\|\mu(X)\|}{\sqrt{\mathbb{E}\|X\|^2 / N}}
\]

monotonically across depth.

This confirms that mean regeneration is intrinsic to the operator,
not merely to embeddings.

\section{Design Objective: Mean-Neutral Nonlinearity}

A nonlinear operator $\phi$ is mean-neutral if

\[
\mathbb{E}[\phi(z)] = 0
\quad \text{whenever } z \text{ is symmetric}.
\]

This holds if and only if $\phi$ is an odd function:

\[
\phi(-z) = -\phi(z).
\]

For example:

\[
\phi(z) = \tanh(z)
\]

satisfies

\[
\mathbb{E}[\tanh(z)] = 0
\quad \text{for symmetric } z.
\]

Thus replacing SwiGLU with an odd nonlinearity
prevents spontaneous mean drift.

However, pure $\tanh$ removes gating and multiplicative structure.

\section{Mean-Neutral Gated Alternatives}

To preserve gating while eliminating mean bias,
we construct symmetric gated activations.

\subsection{Odd Gated Unit}

Define

\begin{equation}
\text{OddGLU}(X)
=
(XW_1) \odot \tanh(XW_2).
\end{equation}

Because $\tanh$ is odd,
if $XW_1$ and $XW_2$ are symmetric,
then the output remains zero-mean.

This preserves multiplicative gating
while eliminating mean regeneration.

\subsection{Centered SwiGLU}

An alternative is explicit centering:

\begin{equation}
A = \text{SwiGLU}(X),
\qquad
A_c = A - \mu(A).
\end{equation}

This enforces zero token-mean regardless of activation type.

This method is architecture-agnostic and preserves expressivity.

\section{Residual Architecture Implications}

In deep Transformers, residual accumulation amplifies any nonzero mean.
Thus even small mean shifts become significant:

\[
\mu(X_L) = \sum_{\ell=1}^{L} \mu(Y_\ell).
\]

Mean-neutral nonlinearities ensure

\[
\mu(Y_\ell) \approx 0,
\]

so the residual stream remains geometrically stable.

Empirical simulations confirm:

\[
R_{\text{none}} \gg R_{\text{pre-block}} \gg R_{\text{post-residual}} = 0.
\]

Post-residual centering is sufficient to prevent mean growth,
independent of activation choice.

\section{Trade-Offs and Practical Considerations}

Replacing SwiGLU with $\tanh$:

\begin{itemize}
\item Eliminates mean regeneration,
\item But reduces expressive gating power.
\end{itemize}

Using OddGLU:

\begin{itemize}
\item Preserves multiplicative interaction,
\item Prevents spontaneous bias growth.
\end{itemize}

Using explicit centering:

\begin{itemize}
\item Works with any activation,
\item Is hardware-friendly,
\item Can be fused into kernels.
\end{itemize}

The optimal choice depends on whether the goal is
full mean-neutrality or preserving next-token likelihood fidelity.

\section{Conclusion}

SwiGLU and related non-odd nonlinearities are a principal source of
mean regeneration in deep Transformer architectures.
This mean growth is not incidental but structural,
arising from the asymmetry of the activation function.

From a geometric perspective, replacing or modifying SwiGLU
is essential if one seeks to design a fully mean-neutral Transformer.
Odd gated units or explicit centering operations
provide principled alternatives.

The broader implication is that Transformer architecture
was not originally designed with high-dimensional mean stability in mind.
By rethinking nonlinear operators from first principles,
we can construct models that maintain semantic contrast
without accumulating harmful bias.